% !TeX root = ../动手学强化学习笔记.tex
% 由主文档通过 \input 载入；请编译主文档。
\setcounter{chapter}{2} % 保留原书章号。
\chapter{多臂赌博机与探索利用}
\label{ch:rl-bandits}

多臂赌博机（multi-armed bandit）保留了强化学习的评价性反馈与探索困难，但暂时移除了动作对未来状态的影响。它使我们能够把“学得准确”与“学习期间获得更多收益”这两个目标分开考察。

\section{赌博机模型与评价指标}

设动作集为 $\mathcal A=\{1,\ldots,K\}$。每个动作 $a$ 对应一个固定奖励分布，均值 $q_*(a)$ 未知；每次选择该动作，便从该分布取得一个新样本。基本随机模型假定各次奖励抽样相互独立，动作选择可以依赖过去数据。记 $A_t$ 为第 $t$ 次选择，反馈仍记为 $R_{t+1}$，则 $\mathbb E[R_{t+1}\mid H_t,A_t=a]=q_*(a)$，其中 $H_t$ 为过去交互历史。

令 $q_*^{\max}=\max_a q_*(a)$，$\Delta_a=q_*^{\max}-q_*(a)$，$N_T(a)=\sum_{t=0}^{T-1}\mathbf1\{A_t=a\}$。相对于始终选择最优臂的策略，期望伪遗憾（expected pseudo-regret）为
\begin{equation}
\mathcal R_T=\mathbb E\!\left[\sum_{t=0}^{T-1}\Delta_{A_t}\right]
=\sum_a\Delta_a\mathbb E[N_T(a)].
\label{eq:rl-bandit-regret}
\end{equation}
它比较期望收益，避免把最优臂一次不走运误当成决策错误。若有两个臂的均值为 $0.6,0.4$，次优臂平均选择 $100$ 次，则其伪遗憾贡献为 $20$，不依赖其奖励实现值恰好如何。

\section{动作价值估计与探索策略}

对每个动作维护访问次数 $N_t(a)$ 和均值估计 $Q_t(a)$。选到 $a$ 并收到 $r$ 后，令 $N_{t+1}(a)=N_t(a)+1$，再用 $1/N_{t+1}(a)$ 为步长更新；其他动作的数据不变。未尝试动作必须有明确的初始化或优先试探规则。

贪心策略只选择当前 $Q_t$ 最大的动作。$\varepsilon$-贪心（epsilon-greedy）以概率 $1-\varepsilon$ 在最大值动作中均匀选择，以概率 $\varepsilon$ 在全部 $K$ 个动作中均匀选择。若最大值唯一，则该动作的总选择概率为 $1-\varepsilon+\varepsilon/K$，每个其他动作的概率为 $\varepsilon/K$。因此进入探索分支仍可能选到贪心动作。

固定 $\varepsilon>0$ 能保持探索，但即使估计已很准确，仍会持续选择次优臂；有正价值差时通常产生线性累计遗憾。令探索率逐渐减小有利于减少后期损失，却必须保留足够探索。乐观初始化暂时鼓励未尝试动作，但效果依赖奖励尺度和估计更新方式，不能自动替代持续探索。

\begin{example}[初期噪声锁定错误动作]
两个伯努利臂的成功概率为 $0.6$ 与 $0.4$。各试一次后，最优臂失败而次优臂成功的概率为 $0.4\times0.4=0.16$。若后续纯贪心且不再重新探索，最优臂的估计一直停在零；有限时间内可能长期使用较差的动作。问题不在均值公式，而在没有获取纠正估计所需的数据。
\end{example}

\section{置信上界与后验采样}

当奖励位于 $[0,1]$，置信上界（upper confidence bound，UCB）使用估计值加不确定性奖励。各臂先试一次，之后一种经典选择规则为
\[
A_t\in\arg\max_a\left\{Q_t(a)+\sqrt{\frac{2\log t}{N_t(a)}}\right\},\qquad t\geq K.
\]
低访问次数产生较大的奖励项，从而推动试探。式中常数对应一种具体 UCB 规则，不是所有奖励分布的通用配置。对独立平稳有界奖励，它可得到次优臂期望访问次数关于 $T$ 的对数级界\cite{sutton-barto}；完整遗憾证明需同时控制各臂在不同访问次数下的均值偏差，这里只说明其构造依据。

汤普森采样（Thompson sampling）用后验不确定性选择动作。伯努利臂成功概率为 $p_a$，采用先验 $p_a\sim\operatorname{Beta}(u_a,v_a)$。观测 $x_a$ 次成功和 $y_a$ 次失败后，后验为 $\operatorname{Beta}(u_a+x_a,v_a+y_a)$。每轮分别从各臂后验抽取一个可能均值，选择抽样值最大的臂。均匀先验下观测三次成功、一次失败，后验为 $\operatorname{Beta}(4,2)$，均值为 $2/3$；选择规则仍是后验抽样，不是每次比较这个均值。

\section{非平稳奖励与上下文赌博机}

若臂均值随时间变化，累计均值会被早期数据拖住。常数步长或滑动窗口能增加近期数据的权重，但追踪更快通常伴随噪声更大。此时式\eqref{eq:rl-bandit-regret}中的固定最优臂未必是恰当基准，评价应说明比较固定策略还是随时间变化的最优选择。

上下文赌博机（contextual bandit）在每次选择前给出上下文 $X_t$，策略依据 $X_t$ 选择动作，奖励均值也依赖上下文。其基本设定不让当前动作改变未来上下文分布。例如在外部独立到来的用户特征下选择推荐项可以使用该模型；若推荐改变用户之后的偏好或是否继续使用系统，则应考虑完整的序贯决策模型。

\section{习题与实验}
\begin{enumerate}
\item 三臂问题中 $Q=(0.2,0.5,0.5)$，取 $\varepsilon=0.3$。按本章规则求三个动作的选择概率，并说明贪心并列的处理为何必须明确。
\item 在两个均值为 $0.7,0.5$ 的臂上，给定期望访问次数 $(800,200)$，计算伪遗憾。再计算两个独立 $\operatorname{Beta}(1,1)$ 后验下各臂被汤普森采样选中的概率。
\item 固定总交互预算，用多个独立种子比较纯贪心、$\varepsilon$-贪心和 UCB；同时保存累计奖励与伪遗憾。随后在中途交换两臂均值，比较样本均值和常数步长的适应过程。
\end{enumerate}





